 than the exponential rule. The exponential taper instead retains
near-field repulsion while rapidly suppressing the far-field tail,
yielding higher task performance and held-out-context generalization
than the matched alternatives.

\subsubsection{External Validation of Remoteness Tapering}
\label{sec:external_tapers}

We finally transfer the same control principle to Hopper and Countdown.
DRPO weights Hopper transitions by their standardized remoteness from
the current Gaussian policy and Countdown responses by their current
mean token surprisal. Within each task, all methods share the dataset,
initialization, labels, optimizer, and training budget. Because these
external tasks do not provide known counterfactual geometry or exact
near-field utility matching, they test transfer rather than repeat the
controlled causal identification.

The external summary in Figure~\ref{fig:taper_control_and_transfer}
pre-registers the transfer comparison, but its cells remain open pending
the formal runs and terminal audit. We therefore make no task-level
ranking claim from this panel. The controlled results establish the
stability--utility mechanism; the external panel tests whether that
principle transfers to continuous offline control and autoregressive
policy optimization.

\subsection{Overall Task Performance}
\label{sec:overall_performance}

Having established the controlled mechanism, we evaluate DRPO as
a complete policy-optimization method on D4RL locomotion and Countdown.
We compare it with published domain-specific baselines and
protocol-matched controls spanning positive-only learning, global
negative scaling, and alternative remoteness tapers. All
protocol-matched methods use the same data and compute budgets. We
report normalized return for D4RL and verifier success for Countdown,
while task-performance failures, support or variance-boundary events,
and NaN/Inf numerical failures remain separate outcomes.

% Required packages:
% \usepackage{booktabs}
% \usepackage{graphicx}

\subsubsection{D4RL Locomotion}
\label{sec:d4rl_performance}

We evaluate the nine D4RL Gym locomotion tasks formed by HalfCheetah,
Hopper, and Walker2d under the medium, medium-replay, and medium-expert
datasets. These datasets vary in behavior quality and coverage, with
medium-replay containing the broadest mixture of historical policies.
Table~\ref{tab:d4rl_performance} compares DRPO with published offline-RL
baselines under the D4RL evaluation protocol. The shared objectives and
taper functions are defined in Appendix~\ref{app:shared_methods};
D4RL-specific score provenance, uncertainty statistics, remoteness
construction, and implementation details are provided in
Appendix~\ref{app:d4rl_methods}.

\begin{table}[t]
\centering
\caption{Normalized returns on the D4RL-9 Gym locomotion benchmark.}
\label{tab:d4rl_performance}

\small
\setlength{\tabcolsep}{3.9pt}
\renewcommand{\arraystretch}{1.07}

\begin{tabular*}{\columnwidth}{@{\extracolsep{\fill}}lcccc}
\toprule
Dataset
&
CQL$^\dagger$
&
TD3+BC$^\dagger$
&
IQL$^\dagger$
&
\textbf{DRPO}
\\
\midrule

halfcheetah-medium
& 44.0 & \textbf{48.3} & \underline{47.4} & 47.3 \\
hopper-medium
& 58.5 & 59.3 & \textbf{66.3} & \underline{63.3} \\
walker2d-medium
& 72.5 & \underline{83.7} & 78.3 & \textbf{85.5} \\

\midrule

halfcheetah-replay
& \textbf{45.5} & \underline{44.6} & 44.2 & 43.3 \\
hopper-replay
& \underline{95.0} & 60.9 & 94.7 & \textbf{95.8} \\
walker2d-replay
& 77.2 & \textbf{81.8} & 73.9 & \underline{78.1} \\

\midrule

halfcheetah-expert
& \textbf{91.6} & 90.7 & 86.7 & \underline{91.1} \\
hopper-expert
& \textbf{105.4} & 98.0 & 91.5 & \underline{101.8} \\
walker2d-expert
& 108.8 & \underline{110.1} & 109.6 & \textbf{110.4} \\

\midrule

\textit{D4RL-9 total}
& \textit{\underline{698.5}}
& \textit{677.4}
& \textit{692.4}
& \textbf{\textit{716.6}}
\\

\bottomrule
\end{tabular*}

\vspace{2pt}

\begin{minipage}{\columnwidth}
\small
\textit{Note.}
Dataset names are shortened for compactness: ``replay'' denotes
``medium-replay'', ``expert'' denotes ``medium-expert'', and the D4RL
``-v2'' suffix is omitted. Results marked by $\dagger$ are reported by
prior work rather than protocol-matched reruns. Rows report task-level
normalized returns, and \textit{D4RL-9 total} sums the nine task means.
Bold and underlined values denote the best and second-best results among
the methods shown in this table. DRPO results are averaged over 10 runs;
standard deviations and seed-level results are reported in
Appendix~\ref{app:d4rl_methods}.
\end{minipage}

\end{table}

The populated candidate table gives DRPO a D4RL-9 total normalized
return of 716.6. This aggregate and the apparent medium-replay advantage
remain provisional until the registered seed-level and terminal audits
close; no causal explanation is inferred from the task-level ranking.

\subsubsection{Countdown}
\label{sec:countdown_performance}

We evaluate DRPO on held-out Countdown puzzles, where a response is
successful only if it forms a valid arithmetic expression that uses the
provided numbers and reaches the target value. Greedy verifier success
is the primary metric, while pass@$k$ measures sampling-level solution
coverage and valid-expression rate distinguishes reasoning failure from
formatting or execution failure.

All compared policy-optimization methods start from the same SFT
checkpoint and use the same frozen response bank, training budget,
validation split, and test puzzles. We compare DRPO with
AsymRE~\citep{arnal2025asymmetric} and
TOPR~\citep{leroux2025tapered} as off-policy baselines.
The shared objectives and taper functions are defined in
Appendix~\ref{app:shared_methods}; Countdown-specific baseline
implementations, response-bank construction, checkpoint selection, and
evaluation details are provided in
Appendix~\ref{app:countdown_methods}.

\begin{table}[t]
\centering
\caption{Performance on held-out Countdown puzzles.}
\label{tab:countdown_performance}

\small
\renewcommand{\arraystretch}{1.10}
\setlength{\tabcolsep}{5.0pt}

\begin{tabular*}{\columnwidth}{@{\extracolsep{\fill}}lccc}
\toprule
Metric
&
AsymRE
&
TOPR
&
\textbf{DRPO}
\\
\midrule

Greedy success (\%)
& --
& --
& --
\\

Pass@$k$ (\%)
& --
& --
& --
\\

Valid expression (\%)
& --
& --
& --
\\

\bottomrule
\end{tabular*}

\vspace{2pt}

\begin{minipage}{\columnwidth}
\small
\textit{Note.}
All methods use the same SFT checkpoint, frozen response bank,
validation split, test puzzles, and training budget. A dash denotes a
result pending the formal run.
\end{minipage}

\end{table}

Table~\ref{tab:countdown_performance} reserves the greedy success,
pass@$k$, and validity fields for the formal run. Because all entries are
pending, we make no Countdown ranking or improvement claim in this
version.

\section{Conclusion}

Repeated off-policy reuse can transform useful negative feedback into a
self-reinforcing geometric effect. Our analysis connects the per-sample reuse
loop to aggregate regimes of stable displacement, drift, and instability.
DRPO interrupts that loop with a detached remoteness taper, yielding Gaussian
ultimate boundedness and polynomial rather than exponential categorical
suppression. Matched diagnostics and controlled interventions isolate policy
geometry as the source of the far-field update and show why selective control
can outperform both removing and globally scaling negative feedback. The main
open limitation is semantic: remoteness alone cannot distinguish obsolete
negative samples from remote samples that remain aligned with task improvement.
Combining geometric control with a reliable utility or alignment signal is a
natural next step.


% Acknowledgements should only appear in the accepted version.


\bibliography{example_paper,missing_references}
\bibliographystyle{icml2026}

%%%%%%%%%%%%%%%%%%%%%%%%%%%%%%%%%%%%%%%%%%%%%%%%%%%%%%%%%%%%%%%%%%%%%%%%%%%%%%%
%%%%%%%%%%%%%%%%%%%%%%%%%%%%%%%%%%%%%%%%%%%%%%%%%%%%%%%%%%%%%%%%%%%%%%%%%%%%%%%
% APPENDIX
%%%%%%%%%%%%%%%%%%%%%%%%%%%%%%%%%%%%%%%%%%%%%%%%%%%%%%%%%%%%%%%%%%%%%%%%%%%%%%%
%%%%%%%%%%%%%%%%%%%%%%%%%%%%%%%%%%%%%%%%%%%%%%%%%%%%%%%%%%%%%%%%%%%%%%%%%%%%%%%
\newpage
\appendix
\onecolumn
\section{Proofs for the Theoretical Results}
\label{app:theory_proofs}

\paragraph{Proof map.}
The score and reuse results are proved in
Sections~\ref{app:distance_strength}--\ref{app:gaussian_reuse_rate}; the
aggregate regimes and policy-family manifestations are proved in
Sections~\ref{app:aggregate_equilibrium} and
\ref{app:policy_family_manifestations}. The selective-tapering bridge,
reference restoration, Gaussian boundedness, control order, exponential tail,
and categorical results appear in Sections~\ref{app:far_field_bridge},
\ref{app:regularized_restoration}, \ref{app:gaussian_drpo_boundedness},
\ref{app:min_order_control}, \ref{app:exp_vanishing_influence}, and
\ref{app:categorical_self_limiting}--\ref{app:categorical_restoring_balance},
respectively.

\subsection{Proof of Proposition~\ref{prop:distance_strength}}
\label{app:distance_strength}

\begin{proof}
Let \(\delta=a-\mu\). For a Gaussian policy with fixed covariance
\(\Sigma\), the negative log-likelihood expands as
\[
    D_\mu(a)
    =
    C_\Sigma
    +
    \frac{1}{2}
    \delta^\top\Sigma^{-1}\delta,
\]
where
\(C_\Sigma=\frac{d}{2}\log(2\pi)+\frac{1}{2}\log|\Sigma|\)
is independent of \(a\) and \(\mu\). This proves
Equation~\eqref{eq:gaussian_remoteness_mahalanobis}.

The mean score function is
\begin{equation}
    \mathbf{s}_\mu(a)
    =
    \nabla_\mu\log\pi_\mu(a)
    =
    \Sigma^{-1}\delta,
\end{equation}
and hence
\begin{equation}
    R_\mu(a)
    =
    \left\|\mathbf{s}_\mu(a)\right\|_2^2
    =
    \delta^\top\Sigma^{-2}\delta.
    \label{eq:proof_gaussian_score_norm}
\end{equation}
Let \(x=\Sigma^{-1/2}\delta\). Then
\begin{equation}
    D_\mu(a)-C_\Sigma
    =
    \frac{1}{2}\|x\|_2^2,
    \qquad
    R_\mu(a)
    =
    x^\top\Sigma^{-1}x.
\end{equation}
Applying the Rayleigh quotient bounds gives
\begin{equation}
\begin{aligned}
    R_\mu(a)
    &\geq
    2\lambda_{\min}(\Sigma^{-1})
    \bigl(D_\mu(a)-C_\Sigma\bigr), \\
    R_\mu(a)
    &\leq
    2\lambda_{\max}(\Sigma^{-1})
    \bigl(D_\mu(a)-C_\Sigma\bigr).
\end{aligned}
\label{eq:proof_gaussian_score_remoteness_bounds}
\end{equation}
Thus, \(R_\mu(a)\) is linearly comparable to
\(D_\mu(a)-C_\Sigma\), with constants determined only by \(\Sigma\).
Since \(D_\mu(a)-C_\Sigma=\frac12 d_{\mathrm{Mah}}^2(a,\mu)\),
the Gaussian mean-score response grows without bound as the standardized
distance grows.

When \(\Sigma=\sigma^2I\), Equation~\eqref{eq:proof_gaussian_score_norm}
gives
\[
    R_\mu(a)
    =
    \frac{1}{\sigma^4}\|\delta\|_2^2,
    \qquad
    D_\mu(a)-C_\Sigma
    =
    \frac{1}{2\sigma^2}\|\delta\|_2^2.
\]
Therefore,
\[
    R_\mu(a)
    =
    \frac{2}{\sigma^2}
    \bigl(D_\mu(a)-C_\Sigma\bigr),
\]
which proves the isotropic exact relation stated in the proposition.

For the categorical policy, let
\begin{equation}
    p
    =
    \pi_z(a)
    =
    e^{-D_z(a)}.
\end{equation}
The score function with respect to the logits is
\begin{equation}
    \mathbf{s}_z(a)
    =
    \nabla_z\log\pi_z(a)
    =
    e_a-\pi_z.
\end{equation}
Its selected-action component is therefore
\begin{equation}
    \left|
        \frac{\partial}{\partial z_a}
        \log\pi_z(a)
    \right|
    =
    1-p
    =
    1-e^{-D_z(a)},
\end{equation}
which proves Equation~\eqref{eq:categorical_selected_score}. Since
\(1-e^{-D}\) is strictly increasing for \(D\geq0\) and converges to
\(1\) as \(D\rightarrow\infty\), the selected-logit response increases
with surprisal but remains bounded.

For completeness, the squared norm of the full categorical score function is
\begin{equation}
\begin{aligned}
    R_z(a)
    &=
    \left\|
        e_a-\pi_z
    \right\|_2^2 \\
    &=
    (1-p)^2
    +
    \sum_{j\neq a}\pi_z(j)^2.
\end{aligned}
\end{equation}
Because the non-selected probabilities sum to \(1-p\), the
Cauchy--Schwarz inequality gives
\begin{equation}
    \sum_{j\neq a}\pi_z(j)^2
    \geq
    \frac{(1-p)^2}{K-1}.
\end{equation}
The same sum is at most \((1-p)^2\). Hence,
\begin{equation}
\begin{aligned}
    \frac{K}{K-1}(1-p)^2
    &\leq
    R_z(a) \\
    &\leq
    2(1-p)^2.
\end{aligned}
\end{equation}
Substituting \(p=e^{-D_z(a)}\) shows that the full categorical logit-score
norm is bounded, completing the proof.
\end{proof}


\subsection{Proof of Theorem~\ref{thm:reuse_dynamics}}
\label{app:reuse_dynamics}

Before proving the generic reuse result, we verify the convexity condition
for the two policy-output coordinates used in the main text. For
fixed-covariance Gaussian means,
\[
D_\mu(a)=C_\Sigma+\frac12(a-\mu)^\top\Sigma^{-1}(a-\mu),
\qquad
\nabla_\mu^2D_\mu(a)=\Sigma^{-1}\succeq0.
\]
For categorical logits \(z\), with \(p=\mathrm{softmax}(z)\),
\[
D_z(a)=\log\sum_j e^{z_j}-z_a,
\qquad
\nabla_z^2D_z(a)=\operatorname{diag}(p)-pp^\top\succeq0,
\]
because
\[
v^\top(\operatorname{diag}(p)-pp^\top)v
=
\operatorname{Var}_{i\sim p}(v_i)\ge0.
\]

\begin{proof}
Let
\[
    f(u)=D(u),
    \qquad
    g_t=\nabla f(u_t),
    \qquad
    R_t=\|g_t\|_2^2.
\]

We first consider a negative coefficient
$\widehat{A}=-c$ with $c>0$. Defining
$\alpha=\eta c$, the update becomes
\[
    u_{t+1}
    =
    u_t+\alpha g_t.
\]
By convexity of $f$,
\begin{align}
    f(u_{t+1})
    &\geq
    f(u_t)
    +
    g_t^\top(u_{t+1}-u_t) \\
    &=
    f(u_t)+\alpha\|g_t\|_2^2.
\end{align}
Hence,
\[
    D_{t+1}
    \geq
    D_t+\eta|\widehat{A}|R_t.
\]

The gradient of a differentiable convex function is monotone, so
\[
    \bigl(g_{t+1}-g_t\bigr)^\top
    \bigl(u_{t+1}-u_t\bigr)
    \geq0.
\]
Using $u_{t+1}-u_t=\alpha g_t$ gives
\[
    g_{t+1}^\top g_t
    \geq
    \|g_t\|_2^2.
\]
By the Cauchy--Schwarz inequality,
\[
    \|g_{t+1}\|_2\|g_t\|_2
    \geq
    g_{t+1}^\top g_t
    \geq
    \|g_t\|_2^2.
\]
Whenever $\|g_t\|_2>0$, division by $\|g_t\|_2$ yields
\[
    \|g_{t+1}\|_2
    \geq
    \|g_t\|_2,
\]
and therefore $R_{t+1}\geq R_t$. The same result is immediate when
$R_t=0$. The inequality is strict whenever the gradient changes
nontrivially along a direction of positive curvature.

We next consider a positive coefficient
$\widehat{A}=c$ with $c>0$. Let
$\alpha=\eta c$, so that
\[
    u_{t+1}
    =
    u_t-\alpha g_t.
\]
Because $f$ has an $L$-Lipschitz gradient, the descent lemma gives
\begin{align}
    f(u_{t+1})
    &\leq
    f(u_t)
    -
    \alpha\|g_t\|_2^2
    +
    \frac{L\alpha^2}{2}\|g_t\|_2^2 \\
    &=
    f(u_t)
    -
    \alpha
    \left(
        1-\frac{L\alpha}{2}
    \right)
    \|g_t\|_2^2.
\end{align}
If $\alpha\leq1/L$, then
\[
    D_{t+1}
    \leq
    D_t-\frac{\alpha}{2}R_t
    =
    D_t-\frac{\eta\widehat{A}}{2}R_t.
\]

It remains to show that the score-function norm does not increase.
For a convex function with an $L$-Lipschitz gradient, the gradient is
$1/L$-cocoercive:
\[
    \bigl(g_t-g_{t+1}\bigr)^\top
    \bigl(u_t-u_{t+1}\bigr)
    \geq
    \frac{1}{L}
    \|g_t-g_{t+1}\|_2^2.
\]
Let $d=g_{t+1}-g_t$. Since
$u_t-u